\Needspace{8\baselineskip}
\Article{256}
The following general observations will make this solution still clearer.

I. If the order \(O\) is derived from a properly primitive order, \(AA\) divides \(D\); but if \(O\) is improperly primitive or derived from an improperly primitive order, \(A\) is even, \(D\) is divisible by \(\frac14 AA\), and the quotient is \(\equiv1\ (\mathrm{mod.}\,4)\). Thus the square of any divisor of \(A\) divides either \(D\) itself or at least \(4D\), and in the latter case the quotient is always \(\equiv1\ (\mathrm{mod.}\,4)\).

II. If \(aa\) divides \(D\), all the values of the expression \(\surd D\ (\mathrm{mod.}\,aa)\) that lie between \(0\) and \(aa-1\) are \(0,a,2a\ldots aa-a\), so \(a\) is the number of forms in \(V\); but only as many of them are properly primitive as there are numbers in
\[
\frac{D}{aa},\quad\frac{D}{aa}-1,\quad\frac{D}{aa}-4\ldots\frac{D}{aa}-(a-1)^2
\]
\SourceLeaf{285}
that have no common divisor with \(a\). When \(a=1\), \(V\) consists of a single form, \((1,0,-D)\), which is always properly primitive. When \(a\) is \(2\) or any power of \(2\), half of those \(a\) numbers are even and half odd, so \(V\) contains \(\frac12 a\) properly primitive forms. When \(a\) is another prime number \(p\) or a power of a prime number \(p\), three cases must be distinguished: all those \(a\) numbers are relatively prime to \(a\), and so all the forms in \(V\) are properly primitive, if \(\frac{D}{aa}\) is not divisible by \(p\) and is also a quadratic nonresidue modulo \(p\); but if \(p\) divides \(\frac{D}{aa}\), \(V\) contains \(\frac{(p-1)a}{p}\) properly primitive forms; finally, if \(\frac{D}{aa}\) is a quadratic residue modulo \(p\) not divisible by \(p\), \(V\) contains \(\frac{(p-2)a}{p}\) properly primitive forms. All this is proved without difficulty. In general, setting \(a=2^\nu p^\pi q^\chi r^\rho\ldots\), where \(p,q,r\) etc. denote distinct odd primes, the number of properly primitive forms in \(V\) is \(NPQR\ldots\), where we must set
\[
\begin{aligned}
N&=1\quad(\text{if }\nu=0)\text{ or }N=2^{\nu-1}\ (\text{if }\nu>0)\\
P&=p^\pi\quad(\text{if }\tfrac{D}{aa}\text{ is a quadratic nonresidue modulo }p)\text{ or}\\
P&=(p-1)p^{\pi-1}\quad(\text{if }\tfrac{D}{aa}\text{ is divisible by }p\text{})\text{ or}\\
P&=(p-2)p^{\pi-1}\quad(\text{if }\tfrac{D}{aa}\text{ is a quadratic residue modulo }p\\ &\qquad\text{and is not divisible by }p\text{})
\end{aligned}
\]
and \(Q,R\) etc. are to be defined from \(q,r\) etc. in the same way as \(P\) from \(p\).

III. If \(aa\) does not divide \(D\), \(\frac{4D}{aa}\) is an integer and \(\equiv1\ (\mathrm{mod.}\,4)\), and the values of the expression \(\surd D\ (\mathrm{mod.}\,aa)\) are \(\frac12 a,\frac32 a,\frac52 a\ldots aa-\frac12 a\). Thus the number of forms in \(V\) is \(a\), but only as many of them are properly primitive as there are numbers in
\[
\frac{D}{aa}-\tfrac14,\quad\frac{D}{aa}-\tfrac94,\quad\frac{D}{aa}-\tfrac{25}{4}\quad\ldots\quad\frac{D}{aa}-(a-\tfrac12)^2
\]
that are relatively prime to \(a\). Whenever \(\frac{4D}{aa}\equiv1\ (\mathrm{mod.}\,8)\), all these numbers are even, so \(V\) contains no properly primitive form; but when \(\frac{4D}{aa}\equiv5\ (\mathrm{mod.}\,8)\). all those numbers are odd, so all forms in \(V\) are properly primitive if \(a\) is \(2\) or a power of \(2\). In general, in this case there are as many properly primitive forms in \(V\) as there are numbers in that list divisible by no odd prime divisor of \(a\). This number is \(NPQR\ldots\) if \(a=2^\nu p^\pi q^\chi r^\rho\ldots\), where we must set \(N=2^\nu\) and derive \(P,Q,R\) etc. from \(p,q,r\) etc. in the same way as in the preceding case.

\SourceLeaf{286}
IV. In this way, then, the numbers of properly primitive forms in \(V,V',V''\) etc. can be determined. For the sum of all these numbers, the following general rule is found without much difficulty: if \(A=2^\nu\mathfrak{A}^\alpha\mathfrak{B}^\beta\mathfrak{C}^\gamma\ldots\), where \(\mathfrak A,\mathfrak B,\mathfrak C\) etc. denote distinct odd primes, the total number of properly primitive forms in \(V,V',V''\) etc. is \(=\frac{A\mathfrak{nabc}\ldots}{2\mathfrak{ABC}\ldots}\), where we must set
\[
\begin{aligned}
\mathfrak n&=1\ (\text{if }\tfrac{4D}{AA}\equiv1,\ \mathrm{mod.}\,8),\text{ or}\\
\mathfrak n&=2\ (\text{if }\tfrac{D}{AA}\text{ is an integer}),\text{ or}\\
\mathfrak n&=3\ (\text{if }\tfrac{4D}{AA}\equiv5,\ \mathrm{mod.}\,8);\text{ further}\\
\mathfrak a&=\mathfrak A\ (\text{if }\mathfrak A\text{ divides }\tfrac{4D}{AA}\text{}),\text{ or}\\
\mathfrak a&=\mathfrak A\pm1\ (\text{if }\mathfrak A\text{ does not divide }\tfrac{4D}{AA}\text{,}\\ &\qquad\text{taking the upper sign}\\
&\qquad\text{or the lower according as }\tfrac{4D}{AA}\text{ is a nonresidue}\\ &\qquad\text{or a quadratic residue modulo }\mathfrak A)
\end{aligned}
\]
and finally derive \(\mathfrak b,\mathfrak c\) etc. from \(\mathfrak B,\mathfrak C\) in the same way as \(\mathfrak a\) from \(\mathfrak A\). To give the proof in greater detail here. is ruled out by the need to be brief.

V. As for the number of classes supplied by the properly primitive forms in \(V,V',V''\) etc., the following three cases must be distinguished.

\textit{First}, when \(D\) is a negative number, each properly primitive form in \(V,V'\) etc. constitutes a separate class; that is, the number of classes sought is expressed by the formula given in the preceding observation, with two exceptions: when \(\frac{4D}{AA}\) is either \(=-4\) or \(=-3\), or equivalently when \(D\) is either \(=-AA\) or \(=-\frac34 AA\). To prove this theorem, clearly we need only show that two distinct forms from \(V,V',V''\) etc. cannot be properly equivalent. Suppose, then, that \((hh,i,k)\), \((h'h',i',k')\) are two distinct properly primitive forms from \(V,V',V''\) etc. belonging to the same class, and that the former is transformed into the latter by the proper substitution \(\alpha,\beta,\gamma,\delta\). We then have the equations
\[
\begin{gathered}\alpha\delta-\beta\gamma=1,\\ hh\alpha\alpha+2i\alpha\gamma+k\gamma\gamma=h'h',\\ hh\alpha\beta+i(\alpha\delta+\beta\gamma)+k\gamma\delta=i'\end{gathered}
\]
From these we easily conclude, \textit{first}, that \(\gamma\) certainly is not \(=0\) (which would imply \(\alpha=\pm1\), \(hh=h'h'\), \(i'\equiv i\ (\mathrm{mod.}\,hh)\), and thus that the proposed forms are identical, contrary to the
\SourceLeaf{287}
hypothesis); \textit{second}, that \(\gamma\) is divisible by the greatest common divisor of the numbers \(h,h'\); (for, setting this divisor \(=r\), it clearly also divides \(2i,2i'\) but is relatively prime to \(k\); moreover, \(rr\) divides \(hhk-h'h'k'=ii-i'i'\), from which we easily deduce that \(r\) also divides \(i-i'\). But we have \(\alpha i'-\beta h'h'=\alpha i+\gamma k\), so \(\gamma k\) and therefore also \(\gamma\) are divisible by \(r\)); \textit{third}, that \((\alpha hh+\gamma i)^2-D\gamma\gamma=hhh'h'\). Setting \(\alpha hh+\gamma i=rp\), \(\gamma=rq\), then, \(p\) and \(q\) are integers, the latter not \(=0\), and \(pp-Dqq=\frac{hhh'h'}{rr}\). But \(\frac{hhh'h'}{rr}\) is the least number divisible by both \(hh\) and \(h'h'\), and therefore divides \(AA\) and so also \(4D\). Thus \(\frac{4Drr}{hhh'h'}\) is a (negative) integer; setting it \(=-e\), we have \(pp-Dqq=-\frac{4D}{e}\), or \(4=(\frac{2rp}{hh'})^2+eqq\). In this equation the term \((\frac{2rp}{hh'})^2\), being a square less than \(4\), must be either \(0\) or \(1\). In the former case \(eqq=4\), and \(D=-(\frac{hh'}{rq})^2\), from which it follows that \(\frac{4D}{AA}\) is a square with a negative sign and therefore certainly not \(\equiv1\ (\mathrm{mod.}\,4)\), so \(O\) is neither improperly primitive nor derived from an improperly primitive order. Thus \(\frac{D}{AA}\) is an integer, from which we easily deduce that \(e\) is divisible by \(4\), \(qq=1\), \(D=-(\frac{hh'}{r})^2\), and that \(\frac{AA}{D}\) too is an integer. We must therefore have \(D=-AA\), or \(\frac{D}{AA}=-1\), which is the first exception. In the latter case \(eqq=3\), so \(e=3\) and \(4D=-3(\frac{hh'}{r})^2\); thus \(3(\frac{hh'}{rA})^2\) is an integer which, since multiplication by the integer square \((\frac{rA}{hh'})^2\) gives \(3\), can only be \(3\). Hence \(4D=-3AA\), or \(D=-\frac34 AA\), which is the second exception. In all the remaining cases, therefore, all properly primitive forms in \(V,V',V''\) etc. belong to distinct classes. --- For the exceptional cases it is enough to state the results of an investigation that is not difficult but must be omitted here to keep things brief. In the former case the properly primitive forms in \(V,V',V''\) etc. always belong to the same class in pairs, in the latter in threes, so the total number of classes sought is, in the former case, half the value of the expression given in the preceding observation, and in the latter one third of it.

\textit{Second}, when \(D\) is a positive square, each properly primitive form in \(V,V',V''\) etc. constitutes a separate class without exception. For suppose \((hh,i,k)\), \((h'h',i',k')\) are two such distinct properly equivalent forms, and that the former is transformed into the latter by the proper substitution \(\alpha,\beta,\gamma,\delta\). All the arguments used in the preceding case that do not assume \(D\) negative clearly hold here too. Thus, with \(p,q,r\)
\SourceLeaf{288}
denoting the same things as there, \(\frac{4Drr}{hhh'h'}\) is again an integer, but now positive rather than negative and also a square. Setting it \(=gg\), we have \((\frac{2rp}{hh'})^2-ggqq=4\), \textit{which is absurd,} since the difference of two squares cannot be \(4\) unless the smaller square is \(0\). The assumption therefore cannot hold.

For the \textit{third} case, however, where \(D\) is a positive nonsquare, we do not yet have a general rule for comparing the number of properly primitive forms in \(V,V',V''\) etc. with the number of distinct classes resulting from them. We can indeed say that the latter either equals the former or is a divisor of it. We have even discovered a special connection between the quotient of these numbers and the smallest values of \(t,u\) satisfying the equation \(tt-Duu=AA\), but describing it here would take too long. Whether that quotient can be found in every case merely by inspecting the numbers \(D,A\) (as in the preceding cases), we cannot say with certainty. Here are some examples, whose number anyone can easily increase. For \(D=13\), \(A=2\), the number of properly primitive forms in \(V\) etc. is \(3\), all equivalent, so they constitute a single class; for \(D=37\), \(A=2\), there are again three properly primitive forms in \(V\) etc., belonging to three distinct classes; for \(D=588\), \(A=7\), there are eight properly primitive forms in \(V\) etc., constituting four classes; for \(D=867\), \(A=17\), \(V\) etc. contain \(18\) properly primitive forms, and just as many for \(D=1445\), \(A=17\), but for the former determinant they fall into two classes, and for the latter into six.

VI. Applying this general theory to the case. where \(O\) is an improperly primitive order, we conclude that the number of classes in this order is to the total number of classes in the properly primitive order as \(1\) is to the number of distinct properly primitive classes constituted by the three forms \((1,0,-D)\), \((4,1,\frac{1-D}{4})\), \((4,3,\frac{9-D}{4})\). Just one class results when \(D\equiv1\ (\mathrm{mod.}\,8)\), because in this case the second and third forms are improperly primitive; but when \(D\equiv5\ (\mathrm{mod.}\,8)\), all three forms are properly primitive and produce as many distinct classes if \(D\) is negative, except for the single case \(D=-3\), in which they constitute one class. Finally, the case where \(D\) is positive (of the form \(8n+5\)) is among those for which a general rule is still lacking. We can nevertheless say that in this case the three forms
\SourceLeaf{289}
belong either to three distinct classes or to one, never to two. For it is easy to see that if the forms \((1,0,-D)\), \((4,1,\frac{1-D}{4})\), \((4,3,\frac{9-D}{4})\) belong respectively to the classes \(K,K',K''\), then \(K+K'=K'\), \(K'+K'=K''\); thus, if \(K\) and \(K'\) are assumed identical, \(K'\) and \(K''\) are identical too. Likewise, if \(K\) and \(K''\) are assumed identical, \(K'\) and \(K''\) are identical too; finally, since \(K'+K''=K\), the assumption that \(K'\) and \(K''\) are identical implies that \(K\) and \(K''\) coincide too. Thus either all three classes \(K,K',K''\) are distinct or all three are identical. \textit{For example,} below \(600\) there are \(75\) numbers of the form \(8n+5\), including \(16\) determinants for which the former case holds, so the number of classes in the properly primitive order is three times that in the improperly primitive order, namely \(37,\allowbreak 101,\allowbreak 141,\allowbreak 189,\allowbreak 197,\allowbreak 269.\ 325,\allowbreak 333,\allowbreak 349,\allowbreak 373,\allowbreak 381,\allowbreak 389,\allowbreak 405,\allowbreak 485,\allowbreak 557,\allowbreak 573\); for the remaining \(59\) the latter case holds, so the number of classes is the same in both orders.

VII. There is hardly any need to observe that the preceding investigation not only allows us to compare the numbers of classes in different orders of the same determinant, but also applies to any different determinants whose ratio is that of squares. Namely, if \(O\) denotes any order of determinant \(dmm\), and \(O'\) an order of determinant \(dm'm'\), \(O\) can be compared with the properly primitive order of determinant \(dmm\), and this with the order derived from the properly primitive order of determinant \(d\), or, what comes to the same thing for the number of classes, with this order itself. The order \(O'\) can be compared with that same order in exactly the same way.
